\chapter{Search and connect}
\label{ch:search}

\section{Finding what to link to}

The introduction left us with a puzzle. The ``rich get richer'' rule explains
why so many networks are scale-free, but as a literal recipe it asks the
impossible: a newcomer would have to know the degree of every node in the
network in order to favor the well-connected ones. In this chapter we solve the
puzzle for an important family of networks --- the \emph{searchable} ones --- by
replacing that impossible global rule with the most natural local act a newcomer
can perform. It \emph{searches} the network, moving from node to node along the
links, and connects to what it happens to find.

The systems to keep in mind are collections of documents that point at one
another. In the \emph{citation network}, each node is a published article and
each link is a reference from one article to another. In the World Wide Web,
each node is a page and each link is a hyperlink. These are \emph{directed}
networks: a reference, like a hyperlink, points one way. An article cites an
older one, not the reverse; a page links to another, which need not link back.
Because direction matters, each node now has two kinds of degree. Its
\emph{in-degree}, written $k^{in}$, counts the links pointing \emph{at} it ---
how many times a paper has been cited, how many pages link to a given page. Its
\emph{out-degree}, written $k^{out}$, counts the links pointing \emph{out} of it
--- how many references a paper makes, how many hyperlinks a page contains. It is
the in-degree that interests us, because in real document networks it is the
in-degree --- the count of incoming attention --- that is broadly distributed,
with a few hugely-cited classics and a long tail of papers cited once or never.

How does a new document come to point at the old ones? Not by consulting a
global popularity ranking. An author writing a new paper finds references by
reading a few papers and following \emph{their} reference lists, and then the
references of those, and so on \citep{price1965}. A web author discovers pages by
surfing from link to link and by typing queries into a search engine
\citep{brin1998}. In every case the information used is strictly local --- a
document one has reached, and the documents it points to. The claim of this
chapter is that this local searching, repeated by every newcomer, adds up to
exactly the ``rich get richer'' behavior that the introduction had to assume by
fiat, and that it brings the power-law degrees, the clustering hierarchy, and
the degree correlations along with it \citep{vazquez2001,vazquez2003}.

\section{Surfing as a random walk}

To turn ``surfing'' into something we can analyze, we model the surfer as a
\emph{random walker} on the network. A random walk is just what it sounds like:
start somewhere, and step to a randomly chosen neighbor, then to a neighbor of
\emph{that} node, and so on. Our surfer does this with a twist. From whatever
page it is on, with probability $q_e$ it follows one of that page's outgoing
hyperlinks, picked at random; with the remaining probability $1-q_e$ it gets
bored, stops following links, and jumps to a completely random page elsewhere.
The single dial $q_e$, a number between $0$ and $1$, sets how patiently the
surfer explores along links before restarting from scratch.

Now ask: in the long run, what fraction of its time does the walker spend at
each page? Call that fraction the page's \emph{visiting probability}, $v_i$ for
node $i$. A page is visited often if many links point to it and those linking
pages are themselves visited often --- a self-referential condition that turns
out to have a clean form. The remarkable thing, worth flagging before any
algebra, is that this visiting probability is \emph{exactly the quantity that
Google was built on}: it is the \emph{PageRank} of the page, the measure of
importance that ranks search results \citep{brin1998}. So the local rule we are
about to write is not a toy --- it is the very process real search engines compute
and then display to the authors who go on to create new links.

\begin{figure}[t]
\centering
\begin{adjustbox}{max width=\linewidth}\begin{tikzpicture}[>={Stealth[round,length=2.4mm]},
  nd/.style={circle,draw=black!65,fill=black!4,minimum size=6.5mm,inner sep=0pt},
  hub/.style={circle,draw=accent,fill=black!30,minimum size=8.5mm,inner sep=0pt},
  ex/.style={->,draw=black!30,line width=0.6pt},
  walk/.style={->,draw=accent,line width=1.4pt},
  jump/.style={->,draw=accent,line width=1.0pt,dashed}]
  \node[nd]  (n1) at (0,0)      {};
  \node[nd]  (n2) at (1.9,1.05) {};
  \node[nd]  (n3) at (1.7,-1.1) {};
  \node[hub] (n4) at (3.7,0.05) {};
  \node[nd]  (n5) at (3.5,-1.7) {};
  \node[nd]  (n6) at (5.6,1.1)  {};
  \node[nd]  (n7) at (5.6,-0.7) {};
  \draw[ex] (n1) -- (n3);
  \draw[ex] (n3) -- (n4);
  \draw[ex] (n3) -- (n5);
  \draw[ex] (n4) -- (n7);
  \draw[ex] (n5) -- (n7);
  \draw[walk] (n1) -- (n2);
  \draw[walk] (n2) -- (n4);
  \draw[walk] (n4) -- (n6);
  \draw[jump] (n6) -- (n5);
  \node[accent,font=\scriptsize] at (0.7,0.8)  {$q_e$};
  \node[accent,font=\scriptsize] at (4.95,-1.55) {restart, $1-q_e$};
\end{tikzpicture}\end{adjustbox}
\caption{\textbf{Surfing as a random walk.} From the page it is on, the searcher
follows one of the outgoing links with probability $q_e$ (solid arrows) or, with
probability $1-q_e$, jumps to a page chosen at random (dashed arrow). A page
reached along many routes --- like the highlighted hub --- is visited most often;
this visiting frequency is the page's PageRank, and it grows in proportion to the
number of links that point at the page.}
\label{fig:walk-surf}
\end{figure}

The picture to keep in mind is Figure~\ref{fig:walk-surf}. We would like to know
how a page's visiting probability depends on its
in-degree. Getting there exactly is hard, so we make the standard physicist's
move of replacing fluctuating quantities by their averages --- a so-called
\emph{mean-field approximation}, ``mean'' as in average. The next box carries
this out; the one sentence to carry away is stated right after it, so the box
can be skipped.

\begin{calculation}{from surfing to the visiting probability}
The visiting probability of node $i$ obeys a balance condition: the walker
arrives at $i$ either by a random jump (probability $1-q_e$, landing on any of
the $n$ nodes equally) or by following a link into $i$ from some node $j$ that
points to it,
\begin{equation}
  v_i = \frac{1-q_e}{n} + q_e \sum_j a_{ij}\, \frac{v_j}{k^{out}_j},
  \label{eq:walk-pagerank}
\end{equation}
where $a_{ij}=1$ when $j$ points to $i$, and the factor $1/k^{out}_j$ accounts
for $j$ spreading its attention over its $k^{out}_j$ outgoing links. This is the
PageRank equation. To simplify it, replace the sum over the specific nodes
pointing to $i$ by an average. A node is reached along an incoming link from a
node of out-degree $k^{out}$ with probability $k^{out} p_{k^{out}}/\langle
k^{out}\rangle$, and that link is then followed with probability $1/k^{out}$; the
two factors of $k^{out}$ cancel, so each incoming link contributes the same
average amount $\Theta = \langle v\rangle/\langle k^{out}\rangle$. Since node $i$
has $k^{in}_i$ incoming links,
\begin{equation}
  v_i = \frac{1-q_e}{n} + q_e\, \Theta\, k^{in}_i.
  \label{eq:walk-vi-mf}
\end{equation}
\end{calculation}

The payoff is Eq.~\eqref{eq:walk-vi-mf}, and it is simple to state in words: a
page's visiting probability grows in a straight line with its in-degree. Double
the number of links pointing at a page and you roughly double how often the
surfer lands on it. That straight-line relationship is the seed from which
everything in this chapter grows.

\section{An effective ``rich get richer''}

Being visited is not the same as being linked to --- you read many pages you do
not cite. So suppose that each time the surfer visits a page it creates a new
incoming link to it only with some probability $q_v$. We can then ask how fast a
page accumulates new links as a function of how many it already has. This rate
is the network-growth analog of preferential attachment, and computing it is the
crux of the chapter. The derivation is short; its conclusion is in the sentence
after the box.

\begin{calculation}{the rate at which a page gains links}
Let $\nu_s$ surfers be active and $\nu_a$ new pages be added per unit time. The
node count $n$ and link count $e$ grow as
\begin{equation}
  \frac{\partial n}{\partial t} = \nu_a, \qquad
  \frac{\partial e}{\partial t} = \nu_s\, q_v\, \langle v \rangle\, n,
  \label{eq:walk-growth}
\end{equation}
because each of the $\langle v\rangle n$ visits in a walk makes a link with
probability $q_v$. Integrating fixes the average degree and hence $\Theta =
\nu_a/(q_v\nu_s n)$. The rate at which a page of in-degree $k^{in}$ gains its
next link is $A(k^{in}) = q_v\, v(k^{in})$; substituting
Eq.~\eqref{eq:walk-vi-mf} and the value of $\Theta$,
\begin{equation}
  A(k^{in}) = \frac{1}{n}\left[\, q_v(1-q_e) + q_e\, \frac{\nu_a}{\nu_s}\, k^{in} \right].
  \label{eq:walk-attach}
\end{equation}
\end{calculation}

Equation~\eqref{eq:walk-attach} is the central result, and it deserves to be
read slowly. The rate at which a page attracts new links has two parts: a
constant piece, the same for every page, which gives even a brand-new page a
chance of being noticed; and a piece \emph{proportional to its current
in-degree}. That second piece is preferential attachment --- the rich getting
richer --- but notice where it came from. The surfer never consulted a global
ranking of pages; it only followed links, one at a time. The bias toward
well-connected pages appeared on its own, for the simple reason that a page with
many incoming links is a page the wandering surfer reaches by many routes. We
have \emph{derived} the rule the introduction had to assume. The same logic
underlies the ``redirection'' trick of \citet{krapivsky2001}, in which a newcomer
points either to a random node or to whatever that node points to --- a one-step
walk that yields the same effective preference.

\section{The degree distribution}

A linear attachment rate is exactly the ingredient needed to produce a power
law. Working out the resulting distribution of in-degrees requires solving a
balance equation for how many pages sit at each in-degree, which we do in the
box. Readers willing to take the outcome on trust can jump to the discussion
after it.

\begin{calculation}{the in-degree distribution and its exponent}
Let $n_{k^{in}}$ be the number of nodes with in-degree $k^{in}$. It changes as
nodes climb from one in-degree to the next by gaining links, with new nodes
injected at in-degree zero:
\begin{equation}
  \frac{\partial n_{k^{in}}}{\partial t}
  = \nu_s\, A_{k^{in}-1}\, n_{k^{in}-1}
  - \nu_s\, A_{k^{in}}\, n_{k^{in}}
  + \nu_a\, \delta_{k^{in},0}\,,
  \label{eq:walk-rate}
\end{equation}
the symbol $\delta_{k^{in},0}$ being $1$ only at $k^{in}=0$. For a network
growing at a constant rate the ratio $\nu_s/\nu_a=\alpha$ is fixed and the shape
of the distribution settles down, $n_{k^{in}}(t)=N p_{k^{in}}$. Solving for that
stationary shape gives (with $\Gamma$ the gamma function, a smooth
generalization of the factorial)
\begin{equation}
  p_{k^{in}} = \frac{1}{1+a}\,
  \frac{\Gamma[a(\gamma-1)+k^{in}]}{\Gamma[a(\gamma-1)]}\;
  \frac{\Gamma[(1+a)(\gamma-1)+1]}{\Gamma[(1+a)(\gamma-1)+k^{in}+1]},
  \label{eq:walk-pk}
\end{equation}
\begin{equation}
  \gamma = 1 + \frac{1}{q_e}, \qquad a = \alpha\, q_v\,(1-q_e),
  \label{eq:walk-gamma}
\end{equation}
which for large in-degree is the power law
\begin{equation}
  p_{k^{in}} \sim (k^{in})^{-\gamma}.
  \label{eq:walk-asym}
\end{equation}
\end{calculation}

Two features of the result, Eq.~\eqref{eq:walk-gamma}, are worth dwelling on.
First, the exponent $\gamma$ is at least $2$ --- exactly the range seen in real
citation and web networks. Second, and more surprising, $\gamma$ depends
\emph{only} on $q_e$, the surfer's patience for following links. It does not
depend on $q_v$, how often readers actually cite what they read, nor on the
growth rate of the network. Whether authors are generous or stingy with
citations changes how many links there are, but not the \emph{shape} of the
distribution; only the manner of \emph{exploring} the network does that. As the
surfer explores ever more deeply ($q_e\to 1$) the exponent drops to its minimum
of $2$, the most hub-dominated case; as the surfer restarts constantly ($q_e\to
0$) the exponent grows without bound and the heavy tail of hubs disappears.

\section{Cliquishness and correlations}

The same walk that sets the degrees also lays down triangles, so it predicts the
clustering hierarchy of the introduction with no extra assumptions. The reason
is simple to picture: whenever the surfer steps from a page to one of that
page's neighbors and links to both, it has closed a triangle. A short
calculation turns this picture into a formula.

\begin{calculation}{how clustering falls with degree}

Counting the links $e_i$ among the neighbors of node $i$ created by the walk 
gives $\partial e_i/\partial t \approx (1+q_e)\,\partial k^{in}_i/\partial 
t$. Integrating this differential equation, assuming $e = 0$ for $k = 1$, we 
obtain $e = (1 + q_e) (k - 1)$. Inserting this into the definition of the 
clustering coefficient,
$\langle c\rangle_k = 2e(k)/[k(k-1)]$, yields
\begin{equation}
  \langle c \rangle_k \approx \frac{2(1+q_e)}{k}.
  \label{eq:walk-ck}
\end{equation}
\end{calculation}

The message of Eq.~\eqref{eq:walk-ck} is that cliquishness is inversely
proportional to degree: a node with twice the degree has roughly half the
clustering. This is precisely the clustering hierarchy of
Eq.~\eqref{eq:intro-hierarchy}, seen in the Internet, the Web, and the protein
network --- and here it falls out for free from growth by searching. The
intuition is that a node gains a triangle every time the surfer steps from it to
a neighbor, so its triangles grow in step with its degree, while the number of
\emph{possible} triangles grows much faster, like the square of the degree;
their ratio, the clustering, therefore falls off as one over the degree.

The degree correlations are governed by the same single dial $q_e$. Numerical
simulation of the model shows that when $q_e$ is small the network looks
essentially uncorrelated --- the average neighbor degree is flat, as in the
baseline of Eq.~\eqref{eq:intro-knn-unc} --- while as $q_e$ grows the average
neighbor degree develops the downward slope of a \emph{disassortative} network,
hubs sitting amid poorly-connected pages. This matches the real Web, whose
measured exponent $\gamma\approx 2.1$ corresponds through
Eq.~\eqref{eq:walk-gamma} to a patient surfer with $q_e>0.5$, squarely in the
disassortative regime. One knob, $q_e$, thus fixes the degree-distribution
exponent, the cliquishness hierarchy, and the sign of the correlations all at
once --- tying together what the introduction presented as three separate facts.

\section{Two ways to search}

It helps to see two concrete versions of the searching rule, because they behave
differently and bracket the range of possibilities. Figure~\ref{fig:walk-models}
contrasts them.

\begin{figure}[t]
\centering
\begin{adjustbox}{max width=\linewidth}\begin{tikzpicture}[>={Stealth[round,length=2.2mm]},
  nd/.style={circle,draw=black!65,fill=black!4,minimum size=5.5mm,inner sep=0pt,font=\scriptsize},
  new/.style={circle,draw=accent,fill=black!30,minimum size=6mm,inner sep=0pt,font=\scriptsize},
  ex/.style={->,draw=black!30,line width=0.6pt},
  add/.style={->,draw=accent,line width=1.2pt}]
  % panel (a): random-walk model
  \begin{scope}
    \node[nd]  (j)  at (0,0)       {$j$};
    \node[nd]  (m1) at (1.5,0.95)  {};
    \node[nd]  (m2) at (1.7,-0.05) {};
    \node[nd]  (m3) at (1.4,-1.05) {};
    \draw[ex] (j)--(m1); \draw[ex] (j)--(m2); \draw[ex] (j)--(m3);
    \node[new] (i) at (-1.7,-0.05) {$i$};
    \draw[add] (i)--(j);
    \draw[add] (i) to[bend left=16] (m1);
    \node[font=\footnotesize] at (0.2,-1.95) {(a) random-walk model};
  \end{scope}
  % panel (b): recursive-search model
  \begin{scope}[xshift=6.6cm]
    \node[nd]  (jj) at (0,0)       {$j$};
    \node[nd]  (q1) at (1.5,0.95)  {};
    \node[nd]  (q2) at (1.7,-0.05) {};
    \node[nd]  (q3) at (1.4,-1.05) {};
    \node[nd]  (p1) at (3.1,1.35)  {};
    \node[nd]  (p2) at (3.2,-0.7)  {};
    \draw[ex] (jj)--(q1); \draw[ex] (jj)--(q2); \draw[ex] (jj)--(q3);
    \draw[ex] (q1)--(p1); \draw[ex] (q2)--(p2);
    \node[new] (ii) at (-1.7,-0.05) {$i$};
    \draw[add] (ii)--(jj);
    \draw[add] (ii) to[bend left=12]  (q1);
    \draw[add] (ii) to[bend left=3]   (q2);
    \draw[add] (ii) to[bend right=16] (q3);
    \draw[add] (ii) to[bend left=20]  (p1);
    \draw[add] (ii) to[bend right=8]  (p2);
    \node[font=\footnotesize] at (0.7,-1.95) {(b) recursive-search model};
  \end{scope}
\end{tikzpicture}\end{adjustbox}
\caption{\textbf{Two ways to search, when a new node $i$ joins.} (a) In the
random-walk model, $i$ links to a randomly chosen node $j$ and then takes a
single step to one of $j$'s neighbors, linking to it as well. (b) In the
recursive-search model, $i$ follows \emph{all} of $j$'s links, and then their
links in turn, linking to every node its widening search reaches. New links are
drawn in bold, pre-existing links thin and gray.}
\label{fig:walk-models}
\end{figure}

\paragraph{The random-walk model.} The simplest version uses a single surfer
that always links what it visits and starts a fresh search from a new node each
time it stops. Stated as a recipe, beginning from a couple of connected nodes,
repeat:
\begin{itemize}
  \item \emph{Add.} Create a new node with one link pointing to a randomly
    chosen existing node.
  \item \emph{Walk.} Having just linked to a node, then with probability $q_e$
    also link to one of that node's neighbors; otherwise go back to the add
    step.
\end{itemize}
%
Here the surfer takes a single step along the network before possibly 
stopping. This model produces the power-law tail in Eq. (\ref{eq:walk-asym}) 
and the clustering law of Eq.~\eqref{eq:walk-ck} (the latter in very good 
agreement with simulation), while the predicted exponent is a slight 
overestimate --- the price of the mean-field approximation in the box above.

\paragraph{The recursive-search model.} A more thorough searcher does not 
stop after one step but follows \emph{all} the links of each node it reaches, 
then all the links of those, exploring an ever-widening territory 
\citep{vazquez2001}. Specifically,
%
\begin{itemize}
  \item \emph{Add.} Create a new node with one link pointing to a randomly
    chosen existing node.
  \item \emph{Walk.} If a link is created to a node in the network then
    with probability $q_e$ a link is also created to each of its nearest
    neighbors. When no link is created go to the add rule.
\end{itemize}
%
This version has a richer behavior --- in fact a genuine \emph{phase transition},
a sudden change of character as the dial $q_e$ is turned, of the kind familiar
from water freezing into ice. The next box records the two cases that can be
solved exactly and the location of the transition; the meaning is summarized
afterward.

\begin{calculation}{the recursive searcher and its phase transition}

For $q_e=0$ only the add step operates, every node is equally likely to be 
linked and $A(k^{in}) = 1/n$ independently of $k^{in}$. In this case the 
steady state solution $n_{k^{in}}(n) = n\, p_{k^{in}}$ of Eq.~\eqref{eq:walk-rate} is
%
\begin{equation}
  p_{k^{in}} = 2^{-(k^{in}+1)}.
  \label{eq:walk-rec0}
\end{equation}
%
For $q_e=1$ the recursion reaches every neighbor, the attachment rate becomes 
exactly linear, $A(k^{in})=(1+k^{in})/n$ --- genuine preferential attachment, 
generated locally --- and the steady state solution $n_{k^{in}}(n) = n\, p_{k^{in}}$ of Eq.~\eqref{eq:walk-rate} is
%
\begin{equation}
  p_{k^{in}} = \frac{1}{(k^{in}+1)(k^{in}+2)} \sim (k^{in})^{-2}.
  \label{eq:walk-rec1}
\end{equation}
%
In between, we don't have an analytical solution. Yet, the average out-degree 
in the steady state is
%
\begin{equation}
  \langle k^{out}\rangle = \frac{1 - \sqrt{1-4 q_e}}{2 q_e},
  \label{eq:walk-rec-dou}
\end{equation}
which is finite only for $q_e \le q_c = 1/4$; above this threshold it diverges.
\end{calculation}

The exact results bracket the behavior: a fast-decaying exponential when the
searcher barely explores, and a clean power law with exponent $2$ when it
explores exhaustively. The new feature is the threshold $q_c$ in
Eq.~\eqref{eq:walk-rec-dou}. Below it the network has a well-defined average
connectivity; above it the average degree runs away and the in-degree
distribution locks onto the power-law exponent $\gamma=2$, no matter the precise
value of $q_e$. Simulations confirm this picture, with the threshold near $q_c
\approx 0.5$.

The striking part is that scale-invariance --- the power law --- is not produced
by carefully tuning $q_e$ to one magic value, but holds across a whole range of
$q_e$ above the threshold. Physicists have a name for systems that sit
automatically at such a critical, scale-free point without fine-tuning:
\emph{self-organized criticality}, the textbook example being a sandpile that,
however you pour sand on it, keeps adjusting itself to the brink of avalanching.
Here the slow addition of new pages and the fast exploration of the existing
network play the roles of slow pouring and fast avalanching. Read this way, the
ubiquity of scale-free citation and web networks is not a delicate coincidence
but the generic, self-organized outcome of building a network by searching it.

\section{Summary}

Searchable networks grow by a rule with nothing global in it: a newcomer
explores the existing network along its links and connects to what it finds.
Modeling the explorer as a random walker --- the very PageRank process that ranks
web pages --- we found that a page's visiting probability rises in a straight line
with its in-degree, which converts the local act of walking into an effective
``rich get richer.'' From that single mechanism follow a power-law in-degree
distribution whose exponent $\gamma=1+1/q_e$ is set purely by how the network is
explored, a clustering hierarchy in which cliquishness falls as one over the
degree, and tunable, generically disassortative degree correlations. Letting the
searcher explore exhaustively turns the smooth picture into a self-organized
critical transition with a robust exponent of $2$. The puzzle of where
preferential attachment ``comes from,'' for this family of networks, is thereby
dissolved: it comes from walking.
